
\section{Definite-side nonvanishing before the comparison induction}
\label{sec:sr-source-definite-nonzero}

We translate \srcid{ThetaLifts.LatticeBasics.NonZero}.
The conclusion \(\operatorname{NonZeroAn}\) concerns a definite
first space with an arbitrary allowable target signature:
\[
 \begin{split}
 p=0&\ \Longrightarrow\
    \exists\Phi\in\operatorname{PolData}(M),\quad I_V(\Phi)\ne0,\\
 q=0&\ \Longrightarrow\
    \exists\Phi'\in\operatorname{PolData}(M')\cap\operatorname{So}(M'),
                          \quad I_{V'}(\Phi')\ne0 .
 \end{split}                                      \tag{N1}
 \label{eq:sr-source-nonzeroAn}
\]
This is a separate theorem from proportionality of the two period
maps. In particular its proof does not assume \(\operatorname{OneConst}\)
at a lower signature.

We retain the exact inputs of \srcid{nonZeroAn\_of}.
The base assertion \(\operatorname{NonZeroBot}\) gives nonzero
polynomial periods on both sides when \(p=q=0\).
The assertions \(\operatorname{TowerA}\) and
\(\operatorname{TowerB}\) supply the adapted rank reductions,
possibly after replacing the upper Cartans and split bases.
\(\operatorname{EvZero}\) supplies a zero-evaluation map whose
image of the adapted upper polynomial space is the full lower
polynomial space.
\(\operatorname{CTBigZero}\) and
\(\operatorname{CTSmallZero}\) give its actual constant-term
identities, with nonzero scalar.
\(\operatorname{LeviEx}\) supplies a Levi lift of every lower
group element, and \(\operatorname{SplitEx}\) supplies uniqueness
of the two-component splitting.

Three further inputs have distinct roles.
\(\operatorname{PolSch}\) places polynomial test vectors in
the adelic Schwartz space.
\(\operatorname{SchToPol}\), in the definite first-group case,
says that if the period vanishes on the prescribed polynomial
space then it vanishes on all Schwartz vectors.
\(\operatorname{ToSo}\), on the definite small side, replaces
a polynomial vector by one in the actual full-rank cutoff
without changing its period.
These inputs are not replaced by an assertion that different
Cartans have the same polynomial subspace.

\begin{srlemma}[The positive definite side, \srcid{big\_of}]
\label{lem:sr-source-big-nonzero}
Assume \(\operatorname{NonZeroBot}\),
\(\operatorname{TowerB}\), \(\operatorname{EvZero}\),
\(\operatorname{CTBigZero}\), \(\operatorname{SchToPol}\),
\(\operatorname{PolSch}\), and \(\operatorname{LeviEx}\).
For every allowable triple with \(p=0\), the large polynomial
period map is nonzero.
\end{srlemma}

\begin{proof}
Induct on \(q\), allowing the small rank \(a'\) and all
corresponding geometric data to vary.
For \(q=0\) use the large part of \(\operatorname{NonZeroBot}\).

For \(q+1\), write \(a'=b+2\); the rank condition ensures
this is possible. The \(B\)-tower gives an adapted upper
triple \(\widetilde G\), a lower triple \(G_1\) of ranks
\((a,b)\), and their planes and zero evaluation \(e\).
The large first signature remains definite, and the lower
small signature parameter is \(q\).
Induction gives a lower polynomial vector \(\Phi_1\)
with nonzero period.
Surjectivity of \(e\) on the adapted polynomial spaces
gives an adapted upper polynomial \(\Psi\) with
\(e\Psi=\Phi_1\).

Suppose every polynomial period for the original upper
Cartan were zero.
The definite \(\operatorname{SchToPol}\) assertion then makes
the upper period zero for every Schwartz vector.
The adapted \(\Psi\) is Schwartz by \(\operatorname{PolSch}\).
Changing to \(\widetilde G\) retained the rational forms,
Darboux coordinates, vacua, Weil action, Haar measure,
and integration domain. Thus the same period function
\(I_V(\Psi)\) is zero, even though its polynomial-space
membership was established with the adapted Cartan.

The zero period has the splitting with both components zero.
Apply \(\operatorname{CTBigZero}\) to this splitting.
For a Levi lift of any \(h_1\in U(W_1)(\mathbb A_F)\),
and with the final target translate equal to \(1\),
its identity reads
\[
                       0=C\,I_{V,1}(e\Psi)(h_1),
                              \qquad C\ne0.
\]
The identity action of \(\omega(1,1)\) has been used here.
Every \(h_1\) has such a lift by \(\operatorname{LeviEx}\).
Hence the lower period of \(\Phi_1\) is identically zero,
contrary to its choice.
\end{proof}

\begin{srlemma}[The negative definite side, \srcid{small\_of}]
\label{lem:sr-source-small-nonzero}
Assume \(\operatorname{NonZeroBot}\),
\(\operatorname{TowerA}\), \(\operatorname{EvZero}\),
\(\operatorname{CTSmallZero}\), \(\operatorname{SchToPol}\),
\(\operatorname{PolSch}\), \(\operatorname{LeviEx}\), and
\(\operatorname{SplitEx}\).
For every allowable triple with \(q=0\), the small polynomial
period map is nonzero, initially without a cutoff requirement
on the chosen vector.
\end{srlemma}

\begin{proof}
Induct on \(p\), allowing \(a\) and all its geometric data to vary.
The case \(p=0\) is the small part of
\(\operatorname{NonZeroBot}\).
In the successor case write \(a=b+2\) and take the
\(A\)-tower. Its lower triple has ranks \((b,a')\)
and parameters \((p-1,0)\).
Choose a lower polynomial vector \(\Phi_1'\)
with nonzero small period, and lift it through the
onto adapted zero evaluation to an upper polynomial
\(\Psi'\).

If the original small polynomial period were identically zero,
\(\operatorname{SchToPol}\) on this definite first space would
make it zero on every Schwartz vector.
As in the preceding proof, this applies to the adapted
polynomial \(\Psi'\), using \(\operatorname{PolSch}\).

The assertion \(\operatorname{CTSmallZero}\) supplies a function
\(f_\alpha\), the splitting
\(\operatorname{Split}_2(I_{V'}(\Psi');f_\alpha,0)\),
and a nonzero \(C\) such that
\[
             f_\alpha(l)=C\,I_{V',1}(e\Psi')(h_1)
\]
for every Levi lift \(l\) of \(h_1\).
Since \(I_{V'}(\Psi')=0\), uniqueness of the two-component
splitting makes \(f_\alpha=0\).
In the source this uniqueness is obtained by applying
\(\operatorname{SplitEx}\) to the zero large test vector
in the same adapted triple, and comparing the splittings
\((f_\alpha,0)\) and \((0,0)\).
Canceling \(C\), and using \(\operatorname{LeviEx}\),
contradicts the nonzero lower period.
\end{proof}

\begin{srproposition}[The original \srcid{nonZeroAn\_of}]
\label{prop:sr-source-nonzeroAn}
Under the inputs of the two preceding lemmas and
\(\operatorname{ToSo}\), assertion (N1) holds.
\end{srproposition}

\begin{proof}
The large assertion is \srcid{big\_of}.
For the small assertion choose the polynomial vector from
\srcid{small\_of}. Since \(q=0\),
\(\operatorname{ToSo}\) gives a polynomial vector in the
actual cutoff with the same period.
That period remains nonzero.
\end{proof}

Thus the exact dependency order in this module is
\[
 \operatorname{NonZeroBot}
 \ \Longrightarrow\
 \begin{cases}
   \text{induction in }q\text{ with }p=0,\\
   \text{induction in }p\text{ with }q=0
 \end{cases}
 \ \Longrightarrow\ \operatorname{NonZeroAn}.
\]
The inputs \(\operatorname{TowerA/B}\), the zero-evaluation
identities, \(\operatorname{SchToPol}\), and
\(\operatorname{ToSo}\) remain explicitly identified until
their separate producer theorems are translated.
This is not the simultaneous proportionality/nonvanishing
induction used in the earlier alternative draft.
